% GENERATED FILE. Edit canonical lessons, then run npm run generate:books.
% canonical-source-hash: fnv1a64:2c55e8cdcf4ce448
% canonical-lessons: SA-C511-kastham, SA-C511-tapayati, SA-C511-uttarayati, SA-C511-pakah, SA-C511-bhrjjati

\chapter{Heat in the Kitchen}
\label{ch:sa-heat-in-the-kitchen}
\hlchaptermodality{511}

\begin{chapteropening}
\textbf{By the end of this chapter:} \emph{I can talk about the fire, the steam and boiling.}

Complete the last lesson of chapter 511: I can talk about the fire, the steam and boiling.
\end{chapteropening}

\section[\texorpdfstring{\sk{काष्ठम्} (kāṣṭham)}{काष्ठम् (kāṣṭham)}]{\sk{काष्ठम्} (kāṣṭham) — firewood; a piece of wood}

\label{lesson:SA-C511-kastham}



\begin{quote}
Before the new one: say the Sanskrit for a pile, then the Sanskrit for a thorough cleaning.
\end{quote}

\subsection*{You'll want to know: \sk{काष्ठम्}}
\textbf{\sk{काष्ठम्}} — \emph{kāṣṭham} — \textquotedblleft{}firewood; a piece of wood\textquotedblright{}.

\begin{quote}
\sk{माता} \sk{चुल्ल्यां} \sk{काष्ठं} \sk{स्थापयति।}
\end{quote}

\begin{quote}
\emph{mātā cullyāṃ kāṣṭhaṃ sthāpayati.}
\end{quote}

\begin{quote}
Mother puts firewood in the stove.
\end{quote}

A neuter noun.

\begin{cousinweb}
\textbf{\sk{काष्ठम्}} (\emph{kāṣṭham}) is the word MW records for a piece of wood, timber, a stick; MW gives no derivation, so its origin is uncertain (MW).
\end{cousinweb}

\begin{grammarlens}[title={What you've built}]
One more word for describing what you lived and felt.
\end{grammarlens}

\subsection*{Your turn}
\begin{itemize}
\raggedright
  \item \emph{Say it:} \emph{kāṣṭham}
  \item \emph{Say it:} \emph{kāṣṭham}, once more
  \item \emph{Say it:} \emph{mātā cullyāṃ kāṣṭhaṃ sthāpayati.} — Mother puts firewood in the stove.
  \item \emph{Recall:} say the Sanskrit for a pile, then the Sanskrit for a thorough cleaning, then say \emph{kāṣṭham} again
\end{itemize}

\begin{tcolorbox}[breakable,colback=teal!4,colframe=teal!35!black,title={Before you move on}]
What does \sk{काष्ठम्} mean? (\textquotedblleft{}Firewood; a piece of wood\textquotedblright{}.) Say it once more.
\end{tcolorbox}
\section[\texorpdfstring{\sk{तापयति} (tāpayati)}{तापयति (tāpayati)}]{\sk{तापयति} (tāpayati) — makes hot; heats up}

\label{lesson:SA-C511-tapayati}



\begin{quote}
Before the new one: say the Sanskrit for a thorough cleaning, then the Sanskrit for firewood.
\end{quote}

\subsection*{You'll want to know: \sk{तापयति}}
\textbf{\sk{तापयति}} — \emph{tāpayati} — \textquotedblleft{}makes hot; heats up\textquotedblright{}.

\begin{quote}
\sk{सा} \sk{स्थाल्यां} \sk{तैलं} \sk{तापयति।}
\end{quote}

\begin{quote}
\emph{sā sthālyāṃ tailaṃ tāpayati.}
\end{quote}

\begin{quote}
She heats oil in the pan.
\end{quote}

Past (she …): \sk{सा} \sk{जलं} \sk{तापितवती}, \emph{sā jalaṃ tāpitavatī}.

\begin{cousinweb}
\textbf{\sk{तापयति}} (\emph{tāpayati}) comes from the causative of the root √tap, \textquotedblleft{}to be hot, to give out heat\textquotedblright{}, which also gives \emph{tapati}: to make a thing hot (MW).
\end{cousinweb}

\begin{grammarlens}[title={What you've built}]
One more word for describing what you lived and felt.
\end{grammarlens}

\subsection*{Your turn}
\begin{itemize}
\raggedright
  \item \emph{Say it:} \emph{tāpayati}
  \item \emph{Say it:} \emph{tāpayati}, once more
  \item \emph{Say it:} \emph{sā sthālyāṃ tailaṃ tāpayati.} — She heats oil in the pan.
  \item \emph{Recall:} say the Sanskrit for a thorough cleaning, then the Sanskrit for firewood, then say \emph{tāpayati} again
\end{itemize}

\begin{tcolorbox}[breakable,colback=teal!4,colframe=teal!35!black,title={Before you move on}]
What does \sk{तापयति} mean? (\textquotedblleft{}Makes hot; heats up\textquotedblright{}.) Say it once more.
\end{tcolorbox}
\section[\texorpdfstring{\sk{उत्तारयति} (uttārayati)}{उत्तारयति (uttārayati)}]{\sk{उत्तारयति} (uttārayati) — takes down; lifts off}

\label{lesson:SA-C511-uttarayati}



\begin{quote}
Before the new one: say the Sanskrit for firewood, then the Sanskrit for makes hot.
\end{quote}

\subsection*{You'll want to know: \sk{उत्तारयति}}
\textbf{\sk{उत्तारयति}} — \emph{uttārayati} — \textquotedblleft{}takes down; lifts off\textquotedblright{}.

\begin{quote}
\sk{माता} \sk{चुल्ल्याः} \sk{पात्रम्} \sk{उत्तारयति।}
\end{quote}

\begin{quote}
\emph{mātā cullyāḥ pātram uttārayati.}
\end{quote}

\begin{quote}
Mother takes the pot off the stove.
\end{quote}

Past (she …): \sk{सा} \sk{पात्रम्} \sk{उत्तारितवती}, \emph{sā pātram uttāritavatī}.

\begin{cousinweb}
\textbf{\sk{उत्तारयति}} (\emph{uttārayati}) is built on \emph{ud}, \textquotedblleft{}up, out\textquotedblright{}, and the causative of the root √tṝ, \textquotedblleft{}to cross over\textquotedblright{}, which also gives \emph{tarati}; MW gives the causative \emph{uttārayati} as to make come out, to take down, to take off (MW).
\end{cousinweb}

\begin{grammarlens}[title={What you've built}]
One more word for describing what you lived and felt.
\end{grammarlens}

\subsection*{Your turn}
\begin{itemize}
\raggedright
  \item \emph{Say it:} \emph{uttārayati}
  \item \emph{Say it:} \emph{uttārayati}, once more
  \item \emph{Say it:} \emph{mātā cullyāḥ pātram uttārayati.} — Mother takes the pot off the stove.
  \item \emph{Recall:} say the Sanskrit for firewood, then the Sanskrit for makes hot, then say \emph{uttārayati} again
\end{itemize}

\begin{tcolorbox}[breakable,colback=teal!4,colframe=teal!35!black,title={Before you move on}]
What does \sk{उत्तारयति} mean? (\textquotedblleft{}Takes down; lifts off\textquotedblright{}.) Say it once more.
\end{tcolorbox}
\section[\texorpdfstring{\sk{पाकः} (pākaḥ)}{पाकः (pākaḥ)}]{\sk{पाकः} (pākaḥ) — cooking}

\label{lesson:SA-C511-pakah}



\begin{quote}
Before the new one: say the Sanskrit for makes hot, then the Sanskrit for takes down.
\end{quote}

\subsection*{You'll want to know: \sk{पाकः}}
\textbf{\sk{पाकः}} — \emph{pākaḥ} — \textquotedblleft{}cooking\textquotedblright{}.

\begin{quote}
\sk{मम} \sk{मातुः} \sk{पाकः} \sk{उत्तमः} \sk{अस्ति।}
\end{quote}

\begin{quote}
\emph{mama mātuḥ pākaḥ uttamaḥ asti.}
\end{quote}

\begin{quote}
My mother's cooking is excellent.
\end{quote}

A masculine noun.

\begin{cousinweb}
\textbf{\sk{पाकः}} (\emph{pākaḥ}) comes from the root √pac, \textquotedblleft{}to cook\textquotedblright{}, which also gives \emph{pacati}; MW gives \emph{pāka} as cooking, baking, roasting, boiling (MW).
\end{cousinweb}

\begin{grammarlens}[title={What you've built}]
One more word for describing what you lived and felt.
\end{grammarlens}

\subsection*{Your turn}
\begin{itemize}
\raggedright
  \item \emph{Say it:} \emph{pākaḥ}
  \item \emph{Say it:} \emph{pākaḥ}, once more
  \item \emph{Say it:} \emph{mama mātuḥ pākaḥ uttamaḥ asti.} — My mother's cooking is excellent.
  \item \emph{Recall:} say the Sanskrit for makes hot, then the Sanskrit for takes down, then say \emph{pākaḥ} again
\end{itemize}

\begin{tcolorbox}[breakable,colback=teal!4,colframe=teal!35!black,title={Before you move on}]
What does \sk{पाकः} mean? (\textquotedblleft{}Cooking\textquotedblright{}.) Say it once more.
\end{tcolorbox}
\section[\texorpdfstring{\sk{भृज्जति} (bhṛjjati)}{भृज्जति (bhṛjjati)}]{\sk{भृज्जति} (bhṛjjati) — roasts}

\label{lesson:SA-C511-bhrjjati}



\begin{quote}
Before the new one: say the Sanskrit for takes down, then the Sanskrit for cooking.
\end{quote}

\subsection*{You'll want to know: \sk{भृज्जति}}
\textbf{\sk{भृज्जति}} — \emph{bhṛjjati} — \textquotedblleft{}roasts\textquotedblright{}.

\begin{quote}
\sk{सः} \sk{चणकान्} \sk{भृज्जति।}
\end{quote}

\begin{quote}
\emph{saḥ caṇakān bhṛjjati.}
\end{quote}

\begin{quote}
He roasts chickpeas.
\end{quote}

Past (he …): \sk{सः} \sk{चणकान्} \sk{भृष्टवान्}, \emph{saḥ caṇakān bhṛṣṭavān}.

\begin{cousinweb}
\textbf{\sk{भृज्जति}} (\emph{bhṛjjati}) comes from the root √bhrajj, \textquotedblleft{}to roast, to parch, to fry\textquotedblright{}, whose present is \emph{bhṛjjati}; its past participle is \emph{bhṛṣṭa}, and MW compares Greek \emph{phrygō} and Latin \emph{frīgere} (MW).
\end{cousinweb}

\begin{grammarlens}[title={What you've built}]
One more word for describing what you lived and felt.
\end{grammarlens}

\subsection*{Your turn}
\begin{itemize}
\raggedright
  \item \emph{Say it:} \emph{bhṛjjati}
  \item \emph{Say it:} \emph{bhṛjjati}, once more
  \item \emph{Say it:} \emph{saḥ caṇakān bhṛjjati.} — He roasts chickpeas.
  \item \emph{Recall:} say the Sanskrit for takes down, then the Sanskrit for cooking, then say \emph{bhṛjjati} again
\end{itemize}

\begin{tcolorbox}[breakable,colback=teal!4,colframe=teal!35!black,title={Before you move on}]
What does \sk{भृज्जति} mean? (\textquotedblleft{}Roasts\textquotedblright{}.) Say it once more.
\end{tcolorbox}
